\documentclass{template}

\begin{document}

\subsection*{Mentorship Proposal}

In this brief document, I would like to outline three key practical topics I am
most interested in for our better understanding of what areas of Mathematics I
should primarily focus on. I decided not to include purely Cryptographic topics
considering our focus on Pure Mathematics.

\textbf{Understanding ecpairing.} Elliptic Curve bilinear pairing (or
\textit{ecpairing} for short) is a cornerstone of many modern cryptography
protocols. For instance, a large class of zk-SNARKs (succinct non-interactive
argument of knowledge) are based on pairing-friendly elliptic
curves~\cite{pari,groth16}. Very roughly, given a cyclic group $\mathbb{G} =
\langle g \rangle$ of (typically prime) order $q$, the user might conceal some
value $\alpha \in \mathbb{Z}/q\mathbb{Z}$ by computing the point $g^{\alpha} \in
\mathbb{G}$. Given three ``encrypted'' points $(u,v,w)=(g^{\alpha}, g^{\beta},
g^{\gamma})$, one can then check linear properties between $\alpha,\beta,\gamma$
without finding them explicitly. For instance, by checking $u^2v^3 = w^5$ we can
verify whether $2\alpha+3\beta=5\gamma$. However, what if we wanted to check
whether $\alpha \cdot \beta = \gamma$? There is no way to do that without any
more advanced algebraic structure. For that reason, cryptographers used prior
research studies from algebraic geometrists and constructed an
\textit{efficient} and \textit{secure} bilinear map $e: \mathbb{G} \times
\mathbb{G} \to \mathbb{G}_T$ with $\mathbb{G}$ being a special group of points
on elliptic curve\footnote{I am actually oversimplifying the explanation since
$e$ in practice is \textit{assymetric} bilinear map where two inputs are
different elliptic curves, one being the twist of another. But I hope I made the
high-level explanation clear enough.}. Such mapping satisfies
$e(g^{\alpha},g^{\beta})=e(g^{\alpha\beta},g)=e(g,g^{\alpha\beta})=e(g,g)^{\alpha\beta}$.
That said, to check whether $\alpha \cdot \beta = \gamma$ one can check whether
$e(u,v) = e(w,g)$. 

Now, the construction of Weil and Tate pairing are highly complicated and
cryptographers even have simplified 150-page documents just to understand
how it is computed~\cite{pairing-for-beginners}. While I have practically
implemented pairing during one of my projects, I still have no idea (from
the Mathematical perspective) about the nature of this object. So I am
highly interested in understanding this and related concepts of elliptic
curves. Additionally, I would like to explore what makes such construction
secure.

\textbf{Complex Algebraic Structures for Cryptography.} One additional 
interesting question that arises in Cryptography is as follows:
\begin{quote}
    Is it possible to construct an \textbf{efficient} and \textbf{secure}
    multilinear maps (at the very least, efficient 3-way multilinear map
    $e(\cdot,\cdot,\cdot)$)?
\end{quote}
While algebraic geometry has natural generalizations of pairings on curves,
they are not efficiently computable (meaning, we cannot practically
implement them). \href{https://www.youtube.com/watch?v=8WDOpzxpnTE}{Dan
Boneh in his awesome lecture ``Pairings in Cryptography''} provides an
awesome overview of pairings and, in particular, discusses why
cryptographers care about more complex multilinear maps.

That said, it is interesting to check which concepts of algebraic geometry might
become useful for cryptographic applications. I should note that this is not
limited to pairings. For instance, I would be glad to read about Hyperelliptic
cryptography~\cite{hyperelliptic}, where instead of elliptic curves,
cryptographers try to come up with more sophisticated curves. In particular, a
\href{http://pustaka.unp.ac.id/file/abstrak_kki/EBOOKS/Kriptografi%20dan%20Ethical%20Hacking%20B.pdf}{``Handbook of Elliptic and Hyperelliptic Curve Cryptography''} 
seems to be an interesting book on this topic (besides, I happen to have a
physical copy in my office).

\textbf{Optimizing Elliptic Curve-based Protocols.} This is also an exciting
topic! One of my commerical projects consisted in optimizing elliptic curve
operations (namely, scalar multiplication and ecpairing), implemented as the
arithmetical circuit over PlonK-based zero-knowledge system. In layman's terms,
the only difference from standard implementation is that in such arithmetical
systems, modular inverse costs the same as modular multiplication. This, in
turn, resulted in numerous nice mathematical optimizations of these operations.
For instance, one paper suggested using Algebraic tori for such systems to make
implementation faster~\cite{r1cs} (using the same technique in the ``regular''
setup is senseless). Additionally, there are numerous non-standard techniques to
optimize operations over certain elliptic curves: for instance, using projective
space or efficient curve endomorphisms. There are even dedicated books on this
topic: see
\href{https://link.springer.com/book/10.1007/b97644}{``Guide to Elliptic Curve Cryptography''}.

That said, it could be nice to explore and to systematize such studies and even
implement them practically.

\subsection*{References}
\renewcommand{\refname}{\vspace{-1em}}
\bibliographystyle{IEEEtran}
\bibliography{refs}

\end{document}